\section{Hyperuniformity on compact doubly homogeneous
  spaces}\label{sec:hyper-comp-doubly}
The main purpose of this paper is to extend the notion of hyperuniformity to
doubly homogeneous manifolds. Thus we first give a concise description of these
spaces and collect the relevant facts that will be used later.
\subsection{Harmonic analysis on compact doubly homogeneous
  spaces}
\label{sec:comp-doubly-homog}
Let us recall the definition of doubly homogeneous spaces. A metric space
$(X,\theta)$ is called doubly homogeneous, if there exists a group $G$ acting
isometrically on $X$ so that
\begin{equation*}
  \forall x_1,x_2,y_1,y_2\in X: \theta(x_1,x_2)=\theta(y_1,y_2)\Rightarrow
  \exists g\in G: y_1=gx_1\wedge y_2=gx_2.
\end{equation*}
The compact doubly homogeneous Riemannian manifolds have been characterised in
\cite{Wang1952:Two_Point_Compact} (see also
\cite{Helgason2000:Groups_geometric_analysis}):
\begin{align*}
  \mathbb{S}^d&\cong \mathrm{SO}(d+1)/\mathrm{SO}(d)\\
  \mathbb{RP}^{d-1}&\cong\mathrm{O}(d)/(\mathrm{O}(d-1)\times \mathrm{O}(1))\\
  \mathbb{CP}^{d-1}&\cong\mathrm{U}(d)/(\mathrm{U}(d-1)\times \mathrm{U}(1))\\
  \mathbb{HP}^{d-1}&\cong\mathrm{Sp}(d)/(\mathrm{Sp}(d-1)\times
  \mathrm{Sp}(1))\\
  \mathbb{OP}^2&\cong\mathrm{F}_4/\mathrm{Spin}(9).
\end{align*}
We also refer to \cite{Anderson_Dostert_Grabner+2023:riesz_green_energy} for
a concise introduction. The case of the sphere has been studied in
\cite{Brauchart_Grabner_Kusner+2020:hyperuniform_point_sets,
  Brauchart_Grabner_Kusner2019:hyperuniform_deterministic}; thus we mainly
focus on the projective spaces $\mathbb{FP}^{d-1}$ in this paper, where we use
$\mathbb{F}$ to represent the underlying scalar ``field'' ($\mathbb{R}$,
$\mathbb{C}$, $\mathbb{H}$, $\mathbb{O}$). The dimension of the space as a real
manifold is then given by
\begin{equation*}
  D=(d-1)\dim_{\mathbb{R}}(\mathbb{F}),
\end{equation*}
where
\begin{equation*}
  \dim_{\mathbb{R}}(\mathbb{C})=2,\quad
  \dim_{\mathbb{R}}(\mathbb{H})=4,\quad \dim_{\mathbb{R}}(\mathbb{O})=8. 
\end{equation*}
Furthermore, we associate the parameters
\begin{equation*}
  \alpha=\frac{d-1}2\dim_{\mathbb{R}}(\mathbb{F})-1\quad\text{and }
  \beta=\frac12\dim_{\mathbb{R}}(\mathbb{F})-1
\end{equation*}
to the spaces $\mathbb{FP}^{d-1}$. We normalise the geodesic metric $\theta$ to
take values in $[0,\frac\pi2]$, and equip the space with the normalised surface
measure $\sigma$.

The measure of the geodesic ball
$B(x,r)=\{y\in\mathbb{FP}^{d-1}\mid\theta(x,y)<r\}$ is then given by
\begin{equation*}
  \sigma(B(x,r))=\int_0^r A(\theta)\,d\theta,
\end{equation*}
where
\begin{equation*}
  A(r)=\frac{2\Gamma(\alpha+\beta+2)}{\Gamma(\alpha+1)\Gamma(\beta+1)}
  \sin(r)^{2\alpha+1}\cos(r)^{2\beta+1}
\end{equation*}
denotes the surface area of the geodesic sphere
$\mathbb{S}(a,r)=\{x\in\mathbb{FP}^{d-1}\mid \theta(a,x)=r\}$. From now on, we
use the notation
\begin{equation*}
  C_{\alpha,\beta}=
  \frac{2\Gamma(\alpha+\beta+2)}{\Gamma(\alpha+1)\Gamma(\beta+1)}
\end{equation*}
for the normalising constant.

The Laplace operator associated to the underlying metric tensor is then given
by
\begin{equation*}
  \Laplace f=-\frac1{A(r)}\frac\partial{\partial r}
  \left(A(r)\frac{\partial f}{\partial r}\right)+\Laplace_{\mathbb{S}(a,r)}f,
\end{equation*}
where $\Laplace_{\mathbb{S}(a,r)}$ denotes the Laplace operator on the geodesic
sphere. In the following we will restrict our attention to zonal functions
depending only on $r$; for these functions $\Laplace_{\mathbb{S}(a,r)}f$
vanishes. Furthermore, we notice that we use the geometers convention that the
Laplacian has non-negative eigenvalues.

The zonal eigenfunctions of $\Laplace$ are then given by the functions
\begin{equation*}
  P_n^{(\alpha,\beta)}(\cos(2\theta(x,a))),
\end{equation*}
the Jacobi-polynomials with parameters $\alpha$ and $\beta$. The corresponding
eigenvalues are given by
\begin{equation*}
  \lambda_n=4n(n+\alpha+\beta+1).
\end{equation*}
The dimension of the corresponding space of eigenfunctions $V_n$ is then given
by
\begin{equation*}
  m_n=\frac{2n+\alpha+\beta+1}{\alpha+\beta+1}
  \frac{(\alpha+\beta+1)_n(\alpha+1)_n}{n!(\beta+1)_n},
\end{equation*}
where $(x)_n=x(x+1)\cdots(x+n-1)$ denotes the rising factorial (Pochhammer
symbol).

\begin{table}[h]
  \renewcommand{\arraystretch}{1.8}
  \begin{tabular}{l| l| l|l|l}
    $X$  & $\alpha$&$\beta$& $\lambda_k$&$m_k = \dim(V_k)$ \\[2mm]\hline
    $\mathbb{RP}^{d-1}$ & $\frac{d-3}2$&$-\frac12$&$2k (2k+d-2)$ &
    $ \frac{4k+d-2}{d-2}\binom{2k+d-3}{d-3}$\\[2mm]\hline
    $\mathbb{CP}^{d-1}$ & $d-2$&$0$&$4k (k+d-1)$ &
    $ \frac{2k+d-1}{d-1}\binom{d+k-2}{d-2}^2$\\[2mm]\hline
    $\mathbb{HP}^{d-1}$&$2d-3$&$1$ & $4k (k+2d-1)$ &
    $ \frac{2k+2d-1}{(2d-1)(2d-2)} \binom{k+2d-2}{2d-2} \binom{k+2d-3}{2d-3}$
    \\[2mm]\hline
    $\mathbb{OP}^{2}$ & $7$&$3$&$4k (k+11)$ &
    $ \frac{2k+11}{1320} \binom{k+10}{7} \binom{k+7}{7}$\\
  \end{tabular}
  \caption{\label{tab:eigen}
    The eigenvalues and dimensions of the eigenspaces of the
    Laplace operator on the projective spaces $\mathbb{FP}^{d-1}$}
\end{table}

The spaces of eigenfunctions $V_n$ are invariant under the group $G$ acting on
$\mathbb{FP}^{d-1}$. The corresponding representation turns out to be
irreducible. Let $Y_{n,k}(x)$ ($k=1,\ldots,m_n$) denote an orthonormal basis of
$V_n$. Then the addition theorem
\begin{equation}\label{eq:addition}
  \sum_{k=1}^{m_n}Y_{n,k}(x)Y_{n,k}(y)=\frac{m_n}{P_n^{(\alpha,\beta)}(1)}
  P_n^{(\alpha,\beta)}(\cos(2\theta(x,y)))
\end{equation}
is an immediate consequence of the irreducibility of $V_n$. Furthermore, we
recall the formula for integrals of zonal functions
\begin{equation}
  \label{eq:zonal}
  \begin{split}
    &\int_{\mathbb{FP}^{d-1}}f(\cos(2\theta(x,y)))\,d\sigma(x)\\
    =
    &C_{\alpha,\beta}\int_0^{\frac\pi2}f(\cos(2\theta))\sin(\theta)^{2\alpha+1}
    \cos(\theta)^{2\beta+1}\,d\theta,
  \end{split}
\end{equation}
see for instance \cite{Helgason1965:radon_transform}.

Every function 
$f(\cos(2t))\in L^2([0,\frac\pi2],\sin(t)^{2\alpha+1}\cos(t)^{2\beta+1}\,dt)$
can be expanded in terms of the orthogonal system
$P_n^{(\alpha,\beta)}(\cos(2t))$:
\begin{equation}\label{eq:fourier}
  f(\cos(2t))=\sum_{n=0}^\infty \widehat{f}(n)P_n^{(\alpha,\beta)}(\cos(2t)),
\end{equation}
where the coefficients are given by
\begin{equation*}
  \widehat{f}(n)=\frac{C_{\alpha,\beta}m_n}{\left(P_n^{(\alpha,\beta)}(1)\right)^2}
  \int\limits_0^{\frac\pi 2}f(\cos(2t))P_n^{(\alpha,\beta)}(\cos(2t))
  \sin(t)^{2\alpha+1}\cos(t)^{2\beta+1}\,dt.
\end{equation*}
The convergence in \eqref{eq:fourier} is \emph{a priori} only in the
$L^2$-sense; for continuous functions $f$ satisfying
\begin{equation}\label{eq:pos-def}
  \sum_{j,k=1}^Nc_j\overline{c_k}f(\cos(2\theta(x_j,x_k)))\geq0
\end{equation}
for all choices of $N$, $x_1,\ldots,x_N\in X$, and
$c_1,\ldots,c_N\in\mathbb{C}$, Mercer's theorem ensures absolute and uniform
convergence of the series \eqref{eq:fourier} (see for instance
\cite{Ferreira_Menegatto2009:eigenvalues_integral_operators}). Functions
satisfying \eqref{eq:pos-def} are called \emph{positive definite}. It follows
from \cite{Bochner1941:Positive_Definite} that positive definite functions are
characterised by the non-negativity of the coefficients $\widehat{f}(n)$.

\subsection{Hyperuniformity on $\mathbb{FP}^{d-1}$}
We use similar ideas as in
\cite{Brauchart_Grabner_Kusner+2020:hyperuniform_point_sets,
  Brauchart_Grabner_Kusner2019:hyperuniform_deterministic}
to define hyperuniform sequences of point sets in the setting of the spaces
$\mathbb{FP}^{d-1}$. As in the original euclidean setting the concept is based
on the \emph{number variance}.
\begin{definition}\label{def-hyper}
  Let $(X_N)_{N\in\mathbb{N}}$ a sequence of point sets in $\mathbb{FP}^{d-1}$
  with $\lim_{N\to\infty}\#X_N=\infty$. The number variance of this sequence is
  then given by
  \begin{equation}
    \label{eq:variance}
    \begin{split}
      &\mathbb{V}(X_N,r)=\mathbb{V}_x\#(X_N\cap B(x,r))\\
      =&
      \int_{\mathbb{FP}^{d-1}}\left(\sum_{y\in X_N}\mathbbm{1}_{B(x,r)}(y)-
        \#X_N\sigma(B(x,r))\right)^2\,d\sigma(x).
    \end{split}
  \end{equation}
\end{definition}
In a probabilistic setting the number variance is the variance of the number of
points in a geodesic ball of radius $r$. The interest in this quantity stems
from the fact that it can be used as a measure for the quality of a point
distribution in the deterministic as well as the probabilistic setting.

\begin{remark}
  The number variance is intimately related to the $L^2$-discrepancy
  \begin{equation*}
    \mathrm{D}_{L^2}(X_N)=\frac1{\#X_N}
    \left(\int_0^{\frac\pi2}\mathbb{V}(X_N,r)\,dr\right)^{\frac12}.
  \end{equation*}
  This notion is a classical measure for the deviation of the discrete
  distribution $\frac1{\#X_N}\sum_{x\in X_N}\delta_x$ from uniform
  distribution. $L^2$-discrepancy has been studied to find lower bounds for this
  deviation; the theory of \emph{irregularities of distribution} has developed
  techniques to derive lower bounds for various notions of discrepancy. For a
  comprehensive introduction and a very good overview over results in this
  context we refer to \cite{Beck_Chen2008:irregularities_distribution}.
\end{remark}
\begin{remark}\label{rem:invariance}
  Stolarsky's celebrated invariance principle
  \cite{Stolarsky1973:sums_distances_sphere} relates the $L^2$-discrepancy of
  a point set $X_N$ on the sphere $\mathbb{S}^d$ to the sum of mutual chordal
  distances
  \begin{equation}\label{eq:invariance}
    \mathrm{D}_{L^2}(X_N)^2+c_d\frac1{(\#X_N)^2}\sum_{x,y\in X_N}\|x-y\|=
    c_d\iint_{\mathbb{S}^d\times\mathbb{S}^d}\|x-y\|\,d\sigma(x)\,d\sigma(y)
  \end{equation}
  for a suitable constant $c_d$. This shows that the $L^2$-discrepancy is
  minimised for sets maximising the sum of mutual distances.

  This invariance principle was generalised to the projective spaces
  $\mathbb{FP}^{d-1}$ by Skriganov
  \cite{Skriganov2020:stolarskys_invariance_principle}. Furthermore, it is
  observed there that the invariance principle is a special case of a much more
  general relation.
\end{remark}

As in
the classical notion of hyperuniformity we will be interested in the dependence
of the number variance on the radius $r$. Three regimes turn out to be of
interest.
\begin{definition} \label{def:hyperuniformity.projective.space}
  A sequence of point sets $(X_N)_{N\in\mathbb{N}}$ is called
  \begin{enumerate}
  \item hyperuniform for large balls, if for every $r\in(0,\frac\pi2)$
    \begin{equation*}
      \mathbb{V}(X_N,r)=o\left(\#X_N\right),
    \end{equation*}
  \item hyperuniform for small balls, if for every sequence of
    $(r_N)_{N\in\mathbb{N}}$ of radii satisfying
    \begin{align*}
      &\lim_{N\to\infty}r_N=0\\
      &\lim_{N\to\infty}\sigma(B(\cdot,r_N))\#X_N=\infty
    \end{align*}
    the following holds
    \begin{equation*}
      \mathbb{V}(X_N,r)=o\left(\#X_N\sigma(B(\cdot,r_N))\right),
    \end{equation*}
  \item hyperuniform for balls at threshold order, if
    \begin{equation*}
      \limsup_{N\to\infty}\mathbb{V}(X_N,r(\#X_N)^{-\frac1D})=\mathcal{O}(r^{D-1})
    \end{equation*}
    for $r\to\infty$.
  \end{enumerate}
\end{definition}
We notice that hyperuniformity at threshold order is obtained by rescaling and
adapting the classical notion to the compact situation. The motivation for the
order $r^{D-1}$ on the right hand side comes from the order of the surface area
of the geodesic sphere; i.i.d. random points would achieve the order
$r^D$. Furthermore, it is clear that a smaller order of magnitude than
$(\#X_N)^{-\frac1D}$ for the radii would not be sensible, because then the
expected number of points in a ball would tend to $0$.  The other two notions
of hyperuniformity have no obvious counterpart in the classical setting. All
three notions have in common that the number variance is required to have
smaller order of magnitude than the variance of i.i.d. random point sets of the
same cardinality.
\subsection{Uniform distribution}\label{sec:uniform-distribution}
Uniform distribution is a notion that formalises the discretisation of
measures.
\begin{definition}
  A sequence of of point sets $(X_N)_{N\in\mathbb{N}}$ in $\mathbb{FP}^{d-1}$
  with $\lim_{N\to\infty}\#X_N=\infty$ is called uniformly distributed, if
  \begin{equation*}
    \lim_{N\to\infty}\frac1{\#X_N}\sum_{x\in
      X_N}\mathbbm{1}_{B(x,r)}=\sigma(B(x,r))
  \end{equation*}
  for all $x\in \mathbb{FP}^{d-1}$ and $r\in[0,\frac\pi2]$.
\end{definition}
For a comprehensive introduction to uniform distribution we refer to
\cite{Drmota_Tichy1997:sequences_discrepancies_applications,
  Kuipers_Niederreiter1974:uniform_distribution_sequences}.

From the general theory (see
\cite{Kuipers_Niederreiter1974:uniform_distribution_sequences}) it is known
that a sequence of point sets is uniformly distributed, if and only if
\begin{equation*}
  \lim_{N\to\infty}\frac1{\#X_N}\sum_{x\in X_N}f(x)=
  \int_{\mathbb{FP}^{d-1}}f(x)\,d\sigma(x)
\end{equation*}
for all continuous functions $f:\mathbb{FP}^{d-1}\to\mathbb{C}$. From this it
follows immediately (using the denseness of the Laplace eigenfunctions and the
addition theorem \eqref{eq:addition}) that $(X_N)_{N\in\mathbb{N}}$ is
uniformly distributed, if and only if
\begin{equation}\label{eq:Weyl}
  \lim_{N\to\infty}\frac1{(\#X_N)^2}\sum_{x,y\in X_N}P_n^{(\alpha,\beta)}
  (\cos(2\theta(x,y)))=0
\end{equation}
for all $n\geq1$.

We now want to relate hyperuniformity to uniform distribution. Since uniform
distribution is a property modelled after the law of large numbers, which holds
especially for i.i.d. random points, and hyperuniformity is modelled as a
property ``better than i.i.d. random points'', we can only expect that
hyperuniformity implies uniform distribution.

In order to show that in all three regimes, we first derive an expression for
the number variance of $X_N$ in terms of Jacobi polynomials. We observe that
$\mathbbm{1}_{B(x,r)}(y)=\mathbbm{1}_{[0,r)}(\theta(x,y))$ is a zonal
function. Thus we can expand it as a series in terms of Jacobi polynomials in
$\cos(2\theta(x,y))$. Using the formula
\begin{multline*}
  \int_0^rP_n^{(\alpha,\beta)}(\cos(2t))\sin(t)^{2\alpha+1}\cos(t)^{2\beta+1}\,dt\\
  =  \frac1{2n}\sin(r)^{2\alpha+2}\cos(r)^{2\beta+2}
  P_{n-1}^{(\alpha+1,\beta+1)}(\cos(2r))
\end{multline*}
for $n\geq1$ (see
\cite{Magnus_Oberhettinger_Soni1966:formulas_theorems_special}), we obtain
\begin{equation*}
  \mathbbm{1}_{B(x,r)}(y)=\sigma(B(x,r))+
  \sum_{n=1}^\infty\frac{C_{\alpha,\beta}m_n}{P_n^{\alpha,\beta}(1)}a_n(r)
  P_n^{(\alpha,\beta)}(\cos(2\theta(x,y))).
\end{equation*}
with
\begin{equation*}
  a_n(r)=\frac1{2nP_n^{\alpha,\beta}(1)}\sin(r)^{2\alpha+2}\cos(r)^{2\beta+2}
  P_{n-1}^{(\alpha+1,\beta+1)}(\cos(2r)).
\end{equation*}

This equation is valid in the $L^2$-sense. Applying Parseval's identity we
obtain
\begin{equation}
  \label{eq:var-Jacobi}
\mathbb{V}(X_n,r)=
\sum_{n=1}^\infty
    \frac{C_{\alpha,\beta}^2m_n}{P_n^{\alpha,\beta}(1)}a_n(r)^2
    \sum_{x,y\in X_N}P_n^{(\alpha,\beta)}(\cos(2\theta(x,y))).
\end{equation}
This series is absolutely and uniformly (as a function of $x$ and $y$)
convergent by Mercer's theorem. Notice that the function
\begin{equation}\label{eq:g_r-formula}
  \begin{split}
    g_r(x,y)&=\!\!\!\!\int\limits_{\mathbb{FP}^{d-1}}\!\!\!
    \left(\mathbbm{1}_{B(x,r)}(z)-\sigma(B(x,r))\right)
    \left(\mathbbm{1}_{B(y,r)}(z)-\sigma(B(y,r))\right)\,d\sigma(z)\\
    &=\sigma(B(x,r)\cap B(y,r))-\sigma(B(\cdot,r))^2
  \end{split}
\end{equation}
is positive definite and that
\begin{equation}\label{eq:v=gr}
  \mathbb{V}(X_N,r)=\sum_{x,y\in X_N}g_r(x,y).
\end{equation}


\subsubsection{Hyperuniformity for large balls implies uniform
  distribution}\label{sec:hyper-large-balls}
Assume now that $(X_N)_{N\in\mathbb{N}}$ is hyperuniform for large balls. For
fixed $n$ choose a distance $0<r<\frac\pi2$ such that
$P_n^{(\alpha,\beta)}(\cos(2r))\neq0$. Then equation \eqref{eq:var-Jacobi}
yields
\begin{multline*}
  0=\lim_{N\to\infty}\frac{\mathbb{V}(X_N,r)}{\#X_N}\geq\\
  \frac{C_{\alpha,\beta}^2m_n}{P_n^{\alpha,\beta}(1)}a_n(r)^2
  \lim_{N\to\infty}\frac1{\#X_N}
  \sum_{x,y\in X_N}P_n^{(\alpha,\beta)}(\cos(2\theta(x,y))).
\end{multline*}
This proves
\begin{theorem}\label{thm:large}
  Let $(X_N)_{N\in\mathbb{N}}$ be hyperuniform for large balls. Then
  $(X_N)_{N\in\mathbb{N}}$ is uniformly distributed. More precisely,
  \begin{equation}\label{eq:Weyl-better}
    \lim_{N\to\infty}\frac1{\#X_N}
  \sum_{x,y\in X_N}P_n^{(\alpha,\beta)}(\cos(2\theta(x,y)))=0
\end{equation}
holds for all $n\geq1$.
\end{theorem}
\begin{remark}
  Notice that \eqref{eq:Weyl-better} gives a better rate of convergence as
  compared to the criterion \eqref{eq:Weyl} for uniform distribution. This
  reflects the heuristic expectation that hyperuniformity should give an
  improved quality of uniform distribution.
\end{remark}

\subsubsection{Hyperuniformity for small balls implies uniform
  distribution}\label{sec:hyper-small-balls}
Assume that $(X_N)_{N\in\mathbb{N}}$ is hyperuniform for small
balls. Furthermore, we observe that for fixed $n\geq1$ the coeffient of
$P_n^{(\alpha,\beta)}$ in \eqref{eq:var-Jacobi} is of order $r_N^{2D}$ for
$N\to\infty$ (notice that $4\alpha+4=2D$). Arguing as above we obtain that
\begin{equation*}
  \lim_{N\to\infty}\frac{\mathbb{V}(X_N,r_N)}{\#X_N\sigma(B(\cdot,r_N))}=0
\end{equation*}
implies
\begin{equation*}
  \lim_{N\to\infty}\frac{r_N^D}{\#X_N}
  \sum_{x,y\in X_N}P_n^{(\alpha,\beta)}(\cos(2\theta(x,y)))=0.
\end{equation*}
Since $r_N$ can be chosen to tend to $0$ arbitrarily slowly, this implies
\begin{equation*}
  \limsup_{N\to\infty}\frac{1}{\#X_N}
  \sum_{x,y\in X_N}P_n^{(\alpha,\beta)}(\cos(2\theta(x,y)))<\infty.
\end{equation*}
Thus we have shown
\begin{theorem}\label{thm:small}
  Let $(X_N)_{N\in\mathbb{N}}$ be hyperuniform for small balls. Then
  $(X_N)_{N\in\mathbb{N}}$ is uniformly distributed. More precisely,
  \begin{equation}\label{eq:Weyl-better2}
    \limsup_{N\to\infty}\frac1{\#X_N}
  \sum_{x,y\in X_N}P_n^{(\alpha,\beta)}(\cos(2\theta(x,y)))<\infty
\end{equation}
holds for all $n\geq1$.
\end{theorem}
\begin{remark}
  Notice that we have an improvement in the quality of uniform distribution in
  comparison with \eqref{eq:Weyl} by a power $(\#X_N)^{1-\eps}$ for $\eps>0$.
\end{remark}
\subsubsection{Hyperuniformity for balls at threshold order implies
  uniform distribution}\label{sec:hyper-balls-threshold}
Assume that $(X_N)_{N\in\mathbb{N}}$ is hyperuniform for balls at threshold
order. For fixed $n\geq 1$ we notice that
\begin{equation*}
  a_n(r(\#X_N)^{-\frac1D})^2
  \sim
  \frac{P_{n-1}^{(\alpha+1,\beta+1)}(1)^2}{(2nP_n^{(\alpha,\beta)}(1))^2}
  \frac{r^{2D}}{(\#X_N)^2}.
\end{equation*}
Combining this with \eqref{eq:var-Jacobi} gives
\begin{multline*}
  r^{2D}\frac{C_{\alpha,\beta}^2m_nP_{n-1}^{(\alpha+1,\beta+1)}(1)^2}
  {(2n)^2(P_n^{\alpha,\beta}(1))^3}
  \limsup_{N\to\infty}\frac1{(\#X_N)^2}
  \sum_{x,y\in X_N}P_n^{(\alpha,\beta)}(\cos(2\theta(x,y)))\\
  \leq
  \limsup_{N\to\infty}\mathbb{V}(X_N,r(\#X_N)^{-\frac1D})=\mathcal{O}(r^{D-1}).
\end{multline*}
This can only hold, if
\begin{equation*}
  \limsup_{N\to\infty}\frac1{(\#X_N)^2}
  \sum_{x,y\in X_N}P_n^{(\alpha,\beta)}(\cos(2\theta(x,y)))=0.
\end{equation*}
Summing up, we have proved
\begin{theorem}
  Let $(X_N)_{N\in\mathbb{N}}$ be hyperuniform for balls at threshold
  order. Then $(X_N)_{N\in\mathbb{N}}$ is uniformly distributed.
\end{theorem}
\begin{remark}
  In this case there seems to be no improvement in the order of convergence of
  the sum \eqref{eq:Weyl}.
\end{remark}
%%% Local Variables:
%%% mode: latex
%%% TeX-master: "Main"
%%% End:
